%LITERATURE REVIEW

\chapter{Literature Survey}

In Chapter 2, students demonstrate the "quality of survey/need identification" and "clarity of problem definition based on need"  as required by the syllabus. This chapter provides the scholarly grounding for the project.

\section{Existing Systems}

This section is an analysis of current solutions based on the survey results students are required to conduct.

Current Landscape: Describe the popular web applications or software tools currently used to address the identified problem.

Methodology of Current Tools: Explain the technologies or Machine Learning models these existing systems typically employ.

Functionality: Detail the features provided by these systems and how they serve the current user base.

\section{Comparative Study}
This section summarizes various methods found in research to "inculcate the process of self-learning and research".

Algorithm Comparison: Create a table comparing different ML algorithms (e.g., SVM vs. Random Forest vs. CNN) used in previous research papers.

Performance Metrics: Compare results from different theoretical or experimental simulations found in the literature, such as accuracy, latency, or resource consumption.

Technical Constraints: Discuss the hardware and software limitations associated with each researched method.

Studies in this field examine various aspects, including data collection, feature extraction, model training, and performance evaluation, with each study presenting unique insights into the detection process. Table \ref{tab:depression_ML} highlights their datasets, methods, and results for quick comparison.

\begin{table}[htbp]
	\centering
	\caption{Summary of depression detection using machine learning techniques}
	\label{tab:depression_ML}
	\small
	\begin{tabular}{
			>{\RaggedRight}p{0.6cm}    % Ref. (narrow)
			>{\RaggedRight}p{2cm}  % Authors (wider)
			>{\RaggedRight}p{0.6cm}  % Year		
			>{\RaggedRight}p{5cm}    % Results
			>{\RaggedRight}p{5cm}    % Limitations
		}
		\toprule
		\textbf{Ref.} & \textbf{Author(s)} & \textbf{Year} & \textbf{Results} & \textbf{Limitations} \\
		\midrule
		\cite{b1} & R. Chiong et al. & 2021 & 98\% accuracy with Gradient Boosting; Combines lexical and content features for robustness, strong feature mix, handles uneven data well, beats baseline results. & Needs labelled data, Does not work for other languages, focuses only on Twitter, uses default settings, manual feature design. \\
		\hline
		\addlinespace
		\cite{b2} & M. S. Liaw et al. & 2022 & 82\% F1-score; User engagement metrics improve detection, better user behavior insights & Data privacy concerns; Network features not significant, Some text incomplete or unclear \\
		\bottomrule
		\hline
	\end{tabular}
\end{table}

\section{Research Gap}

This is the most critical part of the survey, where students identify why a new solution is required to fulfill a societal or research need.

\begin{itemize}
	\item Limitations of Existing Systems: Point out specific shortcomings in current solutions, such as high cost, low accuracy, or lack of a user-friendly web interface.
	\item Unmet Needs: Highlight the specific societal or research requirements that remain unaddressed by the literature surveyed.
	\item Justification for New Project: Clearly state how the proposed "Mini Project - 2 B" will fill these gaps using domain-specific skill sets and engineering fundamentals.	
\end{itemize}


